\documentclass{report}
\usepackage{amsmath,amssymb}
\setlength{\parindent}{0mm}
\setlength{\parskip}{1em}
\begin{document}
\begin{center}
\rule{6in}{1pt} \
{\large Xiaoye, Sherry Li \\
{\bf Enhancing Scalability and Robustness of the Schur Complement Method Using Hierarchical Parallelism}}

Lawrence Berkeley National Laboratory \\ MS 50F-1650 \\ 1 Cyclotron Rd \\ Berkeley \\ CA 94720-8139
\\
{\tt xsli@lbl.gov}\\
Ichitaro Yamazaki\end{center}

Modern numerical simulations give rise to sparse linear systems of
equations that are becoming increasingly difficult to solve using
standard techniques. Matrices that can be directly factorized
are limited in size due to large memory and inter-node communication
requirements. Preconditioned iterative solvers require less memory
and communication, but often require an effective preconditioner, which
is not readily available.
The Schur-complement-based domain decomposition method
has great potential for providing reliable solution of large-scale,
highly-indefinite algebraic equations.

We developed a hybrid direct-iterative linear solver PDSLin,
Parallel Domain-decomposition Schur complement based LINear
solver, which employs techniques from both sparse direct and
iterative methods.
In this method, the global system is first partitioned into smaller
interior subdomain systems, which are connected only through
separators. To compute the solution of the global system,
the unknowns associated with the interior subdomain systems are
first eliminated to form the Schur complement system, which is defined
only on the separators. Since most of the fill occurs in
the Schur complement, the Schur complement system is solved
using a preconditioned iterative method. Then, the solution on the
subdomains is computed by using this solution on the separators and
solving another set of subdomain systems.

For a scalable implementation, it is imperative to exploit multiple
levels of parallelism; namely, solving independent subdomains
in parallel and using multiple processors per subdomain.
This hierarchical parallelism can maintain numerical stability
as well as scalability.
In this multilevel parallelization framework, performance
bottlenecks due to load imbalance and excessive communication
occur at two levels: in an intra-processor group assigned to the
same subdomain and among inter-processor groups assigned to different
subdomains. We developed several techniques to address these issues,
such as taking advantage of the sparsity of right-hand-side columns
during sparse triangular solutions with interfaces, load balancing
sparse matrix-matrix multiplication to form update matrices, and
designing an effective asynchronous point-to-point communication
of the update matrices. We present experimental results to demonstrate
that with the aid of these techniques, our hybrid solver can efficiently
solve very large highly-indefinite linear systems on thousands
of processors.


\end{document}
